
Soit \(f:X'\to X\) une pr\'esentation de Chow~: \(X'\) est projectif sur \(S\)\,; et \(f\) est surjectif. Si l'on admet (a), (b), (c) et (iii) pour \(X'\), alors ils en r\'esultent pour \(X\) en vertu de 4.3 et 3.5. On peut donc supposer \(X\) projectif sur \(S\).

Soit \(H\) ample pour \(X/S\). Si pour un \(s\in S\), \((L^{\otimes m}\otimes H)_s\) est ample, alors \((L^{\otimes m}\otimes H)_t\) est ample pour toute g\'en\'erisation \(t\) de \(s\) (EGA III 4.7.1). Gr\^ace \`a l'\'equivalence de 4.6 (i) et (iv), on a donc \'etablit (a).

Pour (b), (c) et 4.7 (iii), on peut supposer \(S\) int\`egre avec \(s\) point g\'en\'erique dans (b) et (c) (Pour (b) et (c), c'est clair, et pour (iii), cela r\'esulte de 3.4 (i) appliqu\'e au morphisme \(\coprod S_i\to S\), o\`u les \(S_i\) sont les composantes irr\'eductibles de \(S\), munies de la structure r\'eduite induite).

Il existe un sch\'ema int\`egre \(S'\) et un morphisme f.p.p.f. \(S'\to S\), tels que toute composante irr\'eductible \(X'_i\) de \(X_{S'}\), munie de la structure r\'eduite induite soit int\`egre sur \(S'\) (XII 1.3). Pour prouver (b) et (iii), il suffit de le faire pour \(X_{S'}/S'\), (dans (b), par descente\,; dans (iii) par r\'ecurrence noeth\'erienne et 3.4 (i)).

Enfin, consid\'erons le morphisme surjectif canonique \(\coprod X'_i\to X_S\). Alors (b) et (iii) pour \(X_{S'}/S'\) r\'esultent de (c) et (iii) pour les \(X'_i/S'\), en vertu de 4.3, 3.5 et 3.7.

On s'est donc ramen\'e \`a prouver (c) et (iii), en supposant \(S\) int\`egre de point g\'en\'erique \(s\), et \(X\) projectif et int\`egre sur \(S\). Soit \({O}_X(1)\) un faisceau tr\`es ample pour \(X/S\). Alors pour tout \(m,n\), \(\chi(L_t^{\otimes m}(n))\) et \(\chi({O}_{X_t}(n))\) ne d\'ependent pas de \(t\in S\) (EGA III 7.9.4). Donc (c) r\'esulte de 4.6 (i) \(\Longleftrightarrow\) (v), qui implique aussi que la famille des classes des faisceaux inversibles \(\tau\)-\'equivalents \`a \(0\) sur les fibres de \(X/S\) n'a qu'un
